
Code generation models face a multi-objective optimization challenge requiring balance between execution correctness, code quality, and edge case handling. We introduce tri-modal reinforcement learning with dynamic feedback scheduling, a framework that integrates execution feedback, AI-based feedback, and human feedback through phase-specific weight schedules. The approach implements a curriculum over feedback modality rather than task difficulty: Phase 1 (0--30\% training progress) emphasizes execution feedback for correctness foundation, Phase 2 (30--70\%) peaks AI feedback for scalable quality refinement, and Phase 3 (70--100\%) increases human feedback to address edge cases. We validate the mechanism across four sub-hypotheses using HumanEval and MBPP benchmarks (1,038 samples after combining). All 12 gate criteria pass, confirming predicted weight patterns emerge and phase-specific objectives are achieved. Phase 1 execution weight dominates (0.$800\rightarrow 0.714)$, Phase 2 AI weight peaks at 0.545 (50\% progress), and Phase 3 human weight increases (0.$400\rightarrow 0.636)$. Phase 1 achieves $30\times$ faster correctness improvement rate (1.520 vs 0.050 later), Phase 2 improves quality by +0.070 without correctness regression (pass@1 ratio 1.032), and Phase 3 conflict cases resolve to intermediate preferences (median 0.2468 $\in$ [0.1, 0.4], zero collapse). This work validates the training mechanism through proof-of-concept implementation using simulated training trajectories with pretrained models, establishing feedback modality curriculum as a viable and testable RL training strategy for future full-scale evaluation.
